\chapter{Familial Adenomatous Polyposis}
\index{Familial Adenomatous Polyposis}
\label{sec:Familial Adenomatous Polyposis}

\section{Overview}
Familial adenomatous polyposis (FAP) is an autosomal dominant condition caused by germline mutations in the APC gene on chromosome 5q21, characterized by the development of hundreds to thousands of adenomatous polyps throughout the colon and rectum, typically beginning in adolescence. This condition accounts for a significant proportion of cancer-related morbidity and mortality worldwide. Early detection through routine screening and advances in treatment have improved survival rates in recent decades.
\section{Causes}
This condition happens because of APC gene mutation leading to uncontrolled cell growth. Cancer develops through a multistep process involving genetic mutations that drive uncontrolled cell growth and proliferation. These mutations may be inherited or acquired through exposure to carcinogens, radiation, or random errors during cell division. The underlying mechanisms involve complex interactions between genetic predisposition and environmental factors. Disruption of normal physiological processes leads to the characteristic features of the condition. Multiple contributing factors may act synergistically to trigger or worsen the condition. Understanding the specific cause helps guide targeted treatment approaches.
\section{Risk Factors}


\begin{itemize}[noitemsep]
  \item Genetic predisposition and family history
  \item Advancing age
  \item Lifestyle factors (diet, exercise, smoking, alcohol)
  \item Environmental exposures
  \item Underlying medical conditions
  \item Medication use
\end{itemize}
\section{Signs and Symptoms}

\subsection{Common symptoms}
\begin{itemize}[noitemsep]
  \item You may experience the development of polyps in adolescence
  \item Without treatment, CRC develops in virtually 100\% of patients by age 40.
  \item Pain or discomfort in the affected area
  \end{itemize}

\subsection{Emergency warning signs}
\begin{itemize}[noitemsep]
  \item Severe or worsening symptoms of familial adenomatous polyposis require immediate medical evaluation
\end{itemize}
\section{Progression and Stages}
The condition may progress through different stages of severity over time. Early stages often involve mild or subtle symptoms that may be overlooked. Moderate stages involve more noticeable symptoms affecting daily function. Advanced stages may involve significant impairment and complications.
\section{Complications}
\begin{itemize}[noitemsep]
  \item Disease progression leading to worsening symptoms
  \item Secondary infections or injuries
  \item Side effects of treatment
  \item Reduced quality of life
  \item Emotional and psychological impact
\end{itemize}
\section{Diagnosis}
Doctors diagnose this based on colonoscopy showing innumerable polyps and genetic testing. Diagnosis typically involves imaging studies (CT, MRI, PET scans) to identify the location and extent of disease. Biopsy with histopathological examination remains the gold standard for confirming the diagnosis and determining the specific type and grade. Comprehensive medical history and physical examination Laboratory tests to assess organ function and rule out other conditions Imaging studies to visualize affected structures Specialized testing depending on the affected system Consultation with relevant specialists Monitoring of disease progression over time

\begin{itemize}[noitemsep]
  \item Comprehensive medical history and physical examination
  \item Laboratory tests to assess organ function
  \item Imaging studies to visualize affected structures
  \item Specialized testing depending on the affected system
  \item Consultation with relevant specialists
\end{itemize}
\section{Prevention}
\begin{itemize}[noitemsep]
  \item Maintain a healthy lifestyle with balanced diet and exercise
  \item Avoid known risk factors
  \item Attend regular medical check-ups for early detection
  \item Follow recommended screening guidelines
\end{itemize}
\section{Treatment Options}
Treatment focuses on prophylactic proctocolectomy with ileal pouch-anal anastomosis (IPAA) recommended in the late teens to early twenties. Treatment planning is multidisciplinary and depends on the cancer type, stage, and patient factors. Management may include surgery, radiation therapy, systemic therapies (chemotherapy, targeted therapy, immunotherapy), or a combination of these modalities. Medications to manage symptoms and address underlying mechanisms Lifestyle modifications and self-care measures Physical therapy and rehabilitation Surgical intervention when indicated Regular monitoring and follow-up care Patient education and support services

\begin{itemize}[noitemsep]
  \item Medications to manage symptoms and address underlying mechanisms
  \item Lifestyle modifications and self-care measures
  \item Physical therapy and rehabilitation
  \item Surgical intervention when indicated
  \item Regular monitoring and follow-up care
\end{itemize}
\section{Living with It}
\begin{itemize}[noitemsep]
  \item Follow your treatment plan as prescribed
  \item Maintain regular follow-up with your healthcare provider
  \item Make necessary lifestyle adjustments
  \item Seek support from family, friends, or support groups
  \item Stay informed about your condition
\end{itemize}
\section{Outlook}
\begin{itemize}[noitemsep]
  \item Most people recover well with prophylactic surgery
  \item Attenuated FAP (AFAP) has fewer polyps (10--100) and later-onset CRC. Gardner syndrome is a variant with extraintestinal manifestations (osteomas, desmoid tumors, congenital enlargement of retinal pigment epithelium).
\end{itemize}
